\subsection{固定位数4包絡の四次中心元と相対中心}
\label{subsec:order4-explicit-center}
前節の固定包絡に対し $K=\mathbb Q(\gamma)$、$\gamma^4=2$ とし、
中心値 $D$ を記号のまま残す。$U_D=B[t^{-1}]_0/(\mathcal G_D/t^3)$ と書く。
$Z_1,Z_2,Z_3$ を \path{cubic_order4_center_certificate.json} の最初の三つの四次中心元の
$t=1$ での像とする。ここで係数は $K[D]$ に入り、各 $Z_i$ は
「$D$ に依らない活動次数4以下の項」と「$D$ 倍の活動次数1以下の項」の和である。
巨大な自由単語行列や有限個の $D$ での外挿は用いない。

\begin{theorem}[全中心の三生成元と任意の底変更]
\label{thm:order4-relative-center}
$\mathcal U=U_D$ を $K[D]$ 上の族として扱うと、次が成り立つ。
\begin{enumerate}
\item 全ての $D$ に対して $Z(U_D)\cap F_4U_D$ は
$1,Z_1,Z_2,Z_3$ を基底とする。三つの四次中心元は中心全体を生成する。
\item 唯一の関係は
\[
 \mathcal F_D(Z_1,Z_2,Z_3)=\Phi(Z_1,Z_2,Z_3)+Q_D(Z_1,Z_2,Z_3)=0,
 \qquad\deg_Z Q_D\le2,
\]
であり、係数は $K[D]$ に入る。ここで
\begin{align*}
 \Phi={}&Z_3^3+\frac{99-70\gamma^2}{8}Z_2^3
 +\frac{8\gamma^3+16\gamma^2-11\gamma-23}{8}Z_1^3\\
 &-\frac{(\gamma-1)(114\gamma^2-161)}4Z_1Z_2Z_3.
\end{align*}
従って中心族は平坦な三次超曲面族である。
\item 任意の可換 $K[D]$ 代数 $A$ に対し
\[
 Z(\mathcal U)\otimes_{K[D]}A\ \simeq\ Z(\mathcal U\otimes_{K[D]}A)
\]
となる。特に非被約な底変更も含む。
\item 各体上の $U_D$ はこの中心上有限、階数16であり、PI次数は4である。
\end{enumerate}
$Q_D$ の全係数はまだこの定理の検算では求めていない。
また中心の特異性だけで H3 を判定するものではない。
\end{theorem}

\begin{proof}
四次中心元の構成を先に述べる。中心三次商を取る前の包絡で、各次数の
$V\otimes B_{n-1}$ を六二次関係 $R_2\otimes B_{n-2}$ で割る逐次商を作る。
次数0から5の次元は正確に $(1,4,10,20,35,56)$ となる。
次いで次数4、5で $\mathcal G_D B_1,\mathcal G_D B_2$ を割る。
これはそれぞれ $4\times35,10\times56$ の小行列であり、商次元は31、46となる。
四次候補の三交換子の行列は $138\times31$、$D$ に関して一次である。
重みを保つ一次 $D$ ansatz を定数体 $K$ 上で解き、独立な三つの最高四次記号を持つ
中心元 $Z_i$ を得る。各交換子の全係数は $K(D)$ 上の恒等式として零であり、
係数に $D$ の分母がないことも確認する。これは
\path{verify_cubic_order4_center.py} の厳密計算である。

$\operatorname{gr}_{\mathrm{act}}U_D=H=B(E,\mathcal L,\rho)$ で、$\rho$ の位数は4である。
$H$ の中心は、ノルム line 束 $\mathcal L_4$ を
$E\to\bar E=E/\langle\rho\rangle$ に降下させた次数3の line 束の切断環である。
実際 $K(E)[z,z^{-1};\rho]$ 内で中心元の次数は4の倍数に限られ、
次数4では巡回線形化の不変切断となる。有限etale被覆の降下でその次元は3であり、
次数3の line 束は楕円曲線を平面三次へ埋め込むので、中心はこの三切断で生成される。
従って中心の $F_4$ 次元は高々4である。
検算した三つの先頭記号は独立だから、この上界を全 $D$ で達成する。

$K[D]$ を任意の $A$ へ変更しても活動随伴環は $H\otimes_K A$ となる。
最高記号を取ると
$\operatorname{gr}Z(\mathcal U_A)\subset Z(H\otimes_K A)=Z(H)\otimes_K A$
である。最後の等号は $K$ 上の有限次元交換子行列の核が任意の $K$ 代数への
底変更と可換であることによる。
任意の中心元の最高記号を $Z_i$ の最高記号の多項式で表し、その持上げを引くことで
活動次数を下げられる。この帰納法で全中心の生成と、全中心の随伴環の等号を得る。
ここには $A$ の被約性を使っていない。

先頭中心関係 $\Phi$ は次のように独立に確認する。
$y=1$ の elliptic 座標環
\[
 K[x,z]/(z^3+\kappa xz+x^3+1),\qquad
 \kappa=(1+6\gamma^2)/(1-\gamma^2)
\]
内で関係の実際の graph に沿って四次中心元を評価する。
$p=1+\gamma,q=1-\gamma$ に対して graph の第二点は
\[
 (q^2xz+py^2:p^2yz+qx^2:xy-pqz^2)
\]
で与えられる。三ステップ全てで元の多重線型二次関係を直接検算する。
十個の三次中心単項式を上の座標環で簡約すると、$138\times10$ の定数体上の
行列を得て、その核は正確に一次元となり、表示した $\Phi$ を得る。
\path{verify_cubic_order4_leading_center.py} はこの全恒等式を検証する。
有限個の曲線上の点での補間ではない。

$\Phi(Z_i)$ は活動次数12の記号が消えるため中心 $F_{11}$ に入り、
中心の次数は4の倍数なので $F_8$ に入る。次数2以下の中心単項式は先頭で独立だから、
唯一の $Q_D$ を引けば関係を得る。任意の他の関係も最高記号の段階で
$\Phi$ の倍数であり、引き算で次数を下げれば、この一本で生成される。
$Z_3^3$ の係数は1だから、中心は $K[D,Z_1,Z_2]$ 上
$1,Z_3,Z_3^2$ を基底とする自由加群である。
同じ先頭記号の帰納法は全 $A$ で有効なので、この表示が任意の底変更と可換する。

最後に $H$ はその中心上有限で、同じ有限個の加群生成元を持ち上げれば
$U_D$ は中心上有限となる。Hilbert級数
\[
 H_H(s)=\frac{1+s+s^2}{(1-s)^2},\qquad
 H_{Z(H)}(s)=\frac{1+s^4+s^8}{(1-s^4)^2}
\]
の最高係数の比は16であり、整域 $U_D$ の中心上階数も16である。
従って中心商斜体は次数4の中心単純除法環となる。
\end{proof}

\begin{lemma}[48次元sliceと有限停止次数]
\label{lem:order4-fortyeight-slice}
全ての $D,c_1,c_2$ について
$U_D/(Z_1-c_1,Z_2-c_2)$ は次元48である。
その活動濾過は次数8で全体に達し、同次Rees表示では次数8から9への $t$ 乗法が同型となる。
従って次数9までの逐次商から、真の48次元正則表現を得られる。
\end{lemma}
\begin{proof}
$H=S/(g)$ は次元2の AS Gorenstein 環であり、balanced local duality により
augmentation local cohomology は次数2以外で零である\cite{YZ}。
先頭中心の $Z_1,Z_2$ で割ると、$\Phi(0,0,Z_3)=Z_3^3$ から有限次元となる。
従ってこの二中心元のイデアルは augmentation に cofinal であり、
$H$ は $K[\overline Z_1,\overline Z_2]$ 上 depth2の有限graded加群となる。
二変数多項式環上の Auslander--Buchsbaum により graded free である。
その基底の次数分布は
\[
 (1-s^4)^2 H_H(s)=(1+s+s^2)(1+s+s^2+s^3)^2
\]
で、総数48、最高次数8となる。
この基底を $U_D$ へ持ち上げ、最高活動記号を引く帰納法を使うと、
$U_D$ は $K[Z_1,Z_2]$ 上もこの基底を持つ自由加群となる。
独立性も先頭記号の独立性から従う。
特殊化後のフィルター基底とRees環での主張が全て従い、
次数9以降に新たな隠れた関係が現れて次元が落ちる可能性も除かれる。
\end{proof}

\begin{remark}[全三次関係の有限補間に必要な事前次数境界]
\label{rem:order4-center-interpolation}
$Q_D$ の三次元に関する次数が $j$ の係数は、$D$ について
\[
 j=2:\ \deg_D\le1,\qquad
 j=1:\ \deg_D\le2,\qquad
 j=0:\ \deg_D\le4
\]
となる。これを標本から推測するのではない。
実際 $a\mapsto r a,\beta\mapsto r^2\beta$ とした二変数族を取り、
$X,r,D,Z_i$ の重みを $1,1,3,4$ とする。
$r^4Z_i(X/r;D/r^3)$ は、構成した四次中心元の $D$ 一次項が活動次数1以下であることから
$r$ の負冪を持たず、$r=0$ にも延長する。
その先頭記号は同じ三生成元なので、上の中心生成・唯一の関係の証明は
$K[r,D]$ 上でも成立する。唯一の monic 関係は重み12であり、表示した境界を得る。

$\Phi$ が既知なので、残る係数空間の次元は26である。
$Z_3$ に関する係数ごとに、その次元は2、7、17となる。
したがって事前に階数を確認した17個の $(D,Z_1,Z_2)$ 評価で全係数を決定できる。
各評価では中心二元で割った48次元代数を次数ごとの商で構成し、
$Z_3$ の monic 三次最小多項式を全行列恒等式として確認する。
補間後には17点の評価空間の可逆性とこの次数境界を使って全パラメータの結論へ移る。
これらの全評価が完了するまでは、全三次式を計算済みと報告しない。
\end{remark}
